% Options for packages loaded elsewhere
\PassOptionsToPackage{unicode}{hyperref}
\PassOptionsToPackage{hyphens}{url}
\PassOptionsToPackage{dvipsnames,svgnames,x11names}{xcolor}
%
\documentclass[
]{article}
\usepackage{amsmath,amssymb}
\usepackage{iftex}
\ifPDFTeX
  \usepackage[T1]{fontenc}
  \usepackage[utf8]{inputenc}
  \usepackage{textcomp} % provide euro and other symbols
\else % if luatex or xetex
  \usepackage{unicode-math} % this also loads fontspec
  \defaultfontfeatures{Scale=MatchLowercase}
  \defaultfontfeatures[\rmfamily]{Ligatures=TeX,Scale=1}
\fi
\usepackage{lmodern}
\ifPDFTeX\else
  % xetex/luatex font selection
\fi
% Use upquote if available, for straight quotes in verbatim environments
\IfFileExists{upquote.sty}{\usepackage{upquote}}{}
\IfFileExists{microtype.sty}{% use microtype if available
  \usepackage[]{microtype}
  \UseMicrotypeSet[protrusion]{basicmath} % disable protrusion for tt fonts
}{}
\makeatletter
\@ifundefined{KOMAClassName}{% if non-KOMA class
  \IfFileExists{parskip.sty}{%
    \usepackage{parskip}
  }{% else
    \setlength{\parindent}{0pt}
    \setlength{\parskip}{6pt plus 2pt minus 1pt}}
}{% if KOMA class
  \KOMAoptions{parskip=half}}
\makeatother
\usepackage{xcolor}
\usepackage[margin=1in]{geometry}
\usepackage{color}
\usepackage{fancyvrb}
\newcommand{\VerbBar}{|}
\newcommand{\VERB}{\Verb[commandchars=\\\{\}]}
\DefineVerbatimEnvironment{Highlighting}{Verbatim}{commandchars=\\\{\}}
% Add ',fontsize=\small' for more characters per line
\usepackage{framed}
\definecolor{shadecolor}{RGB}{248,248,248}
\newenvironment{Shaded}{\begin{snugshade}}{\end{snugshade}}
\newcommand{\AlertTok}[1]{\textcolor[rgb]{0.94,0.16,0.16}{#1}}
\newcommand{\AnnotationTok}[1]{\textcolor[rgb]{0.56,0.35,0.01}{\textbf{\textit{#1}}}}
\newcommand{\AttributeTok}[1]{\textcolor[rgb]{0.13,0.29,0.53}{#1}}
\newcommand{\BaseNTok}[1]{\textcolor[rgb]{0.00,0.00,0.81}{#1}}
\newcommand{\BuiltInTok}[1]{#1}
\newcommand{\CharTok}[1]{\textcolor[rgb]{0.31,0.60,0.02}{#1}}
\newcommand{\CommentTok}[1]{\textcolor[rgb]{0.56,0.35,0.01}{\textit{#1}}}
\newcommand{\CommentVarTok}[1]{\textcolor[rgb]{0.56,0.35,0.01}{\textbf{\textit{#1}}}}
\newcommand{\ConstantTok}[1]{\textcolor[rgb]{0.56,0.35,0.01}{#1}}
\newcommand{\ControlFlowTok}[1]{\textcolor[rgb]{0.13,0.29,0.53}{\textbf{#1}}}
\newcommand{\DataTypeTok}[1]{\textcolor[rgb]{0.13,0.29,0.53}{#1}}
\newcommand{\DecValTok}[1]{\textcolor[rgb]{0.00,0.00,0.81}{#1}}
\newcommand{\DocumentationTok}[1]{\textcolor[rgb]{0.56,0.35,0.01}{\textbf{\textit{#1}}}}
\newcommand{\ErrorTok}[1]{\textcolor[rgb]{0.64,0.00,0.00}{\textbf{#1}}}
\newcommand{\ExtensionTok}[1]{#1}
\newcommand{\FloatTok}[1]{\textcolor[rgb]{0.00,0.00,0.81}{#1}}
\newcommand{\FunctionTok}[1]{\textcolor[rgb]{0.13,0.29,0.53}{\textbf{#1}}}
\newcommand{\ImportTok}[1]{#1}
\newcommand{\InformationTok}[1]{\textcolor[rgb]{0.56,0.35,0.01}{\textbf{\textit{#1}}}}
\newcommand{\KeywordTok}[1]{\textcolor[rgb]{0.13,0.29,0.53}{\textbf{#1}}}
\newcommand{\NormalTok}[1]{#1}
\newcommand{\OperatorTok}[1]{\textcolor[rgb]{0.81,0.36,0.00}{\textbf{#1}}}
\newcommand{\OtherTok}[1]{\textcolor[rgb]{0.56,0.35,0.01}{#1}}
\newcommand{\PreprocessorTok}[1]{\textcolor[rgb]{0.56,0.35,0.01}{\textit{#1}}}
\newcommand{\RegionMarkerTok}[1]{#1}
\newcommand{\SpecialCharTok}[1]{\textcolor[rgb]{0.81,0.36,0.00}{\textbf{#1}}}
\newcommand{\SpecialStringTok}[1]{\textcolor[rgb]{0.31,0.60,0.02}{#1}}
\newcommand{\StringTok}[1]{\textcolor[rgb]{0.31,0.60,0.02}{#1}}
\newcommand{\VariableTok}[1]{\textcolor[rgb]{0.00,0.00,0.00}{#1}}
\newcommand{\VerbatimStringTok}[1]{\textcolor[rgb]{0.31,0.60,0.02}{#1}}
\newcommand{\WarningTok}[1]{\textcolor[rgb]{0.56,0.35,0.01}{\textbf{\textit{#1}}}}
\usepackage{graphicx}
\makeatletter
\def\maxwidth{\ifdim\Gin@nat@width>\linewidth\linewidth\else\Gin@nat@width\fi}
\def\maxheight{\ifdim\Gin@nat@height>\textheight\textheight\else\Gin@nat@height\fi}
\makeatother
% Scale images if necessary, so that they will not overflow the page
% margins by default, and it is still possible to overwrite the defaults
% using explicit options in \includegraphics[width, height, ...]{}
\setkeys{Gin}{width=\maxwidth,height=\maxheight,keepaspectratio}
% Set default figure placement to htbp
\makeatletter
\def\fps@figure{htbp}
\makeatother
\setlength{\emergencystretch}{3em} % prevent overfull lines
\providecommand{\tightlist}{%
  \setlength{\itemsep}{0pt}\setlength{\parskip}{0pt}}
\setcounter{secnumdepth}{-\maxdimen} % remove section numbering
\ifLuaTeX
  \usepackage{selnolig}  % disable illegal ligatures
\fi
\IfFileExists{bookmark.sty}{\usepackage{bookmark}}{\usepackage{hyperref}}
\IfFileExists{xurl.sty}{\usepackage{xurl}}{} % add URL line breaks if available
\urlstyle{same}
\hypersetup{
  pdftitle={Lab 04},
  colorlinks=true,
  linkcolor={Maroon},
  filecolor={Maroon},
  citecolor={Blue},
  urlcolor={Blue},
  pdfcreator={LaTeX via pandoc}}

\title{Lab 04}
\author{}
\date{\vspace{-2.5em}}

\begin{document}
\maketitle

\hypertarget{lab-04-raster-data-management-and-analysis}{%
\section{Lab 04: Raster data management and
analysis}\label{lab-04-raster-data-management-and-analysis}}

\hypertarget{read-the-instructions-completely-before-starting-the-lab}{%
\subsubsection{Read the instructions COMPLETELY before starting the
lab}\label{read-the-instructions-completely-before-starting-the-lab}}

This lab builds directly on the raster analysis concepts we worked
through in class, including the Week 12 in-class notebook and the
Wildcard Thursday activity. You should have that code available for
reference.

\hypertarget{formatting-your-submission}{%
\subsubsection{Formatting your
submission}\label{formatting-your-submission}}

This lab must be placed into a public repository on GitHub
(www.github.com). Before the due date, submit \textbf{on Canvas} a link
to the repository. I will then download your repositories and run your
code. The code must be contained in either a .R script or a .Rmd
markdown document. As I need to run your code, any data you use in the
lab must be referenced using \textbf{relative path names}. Finally,
answers to questions I pose in this document must also be in the
repository at the time you submit your link to Canvas. They can be in a
separate text file, or if you decide to use an RMarkdown document, you
can answer them directly in the doc.

\hypertarget{data}{%
\subsection{Data}\label{data}}

All data for this lab are in the course repository at:

\begin{verbatim}
../data/wildcard_thursday/erie_raster_ex/
\end{verbatim}

You will find two subdirectories:

\begin{itemize}
\tightlist
\item
  \texttt{images/} --- seven weekly composite GeoTIFF rasters of
  cyanobacteria bloom intensity in Lake Erie (summer--fall 2017)
\item
  \texttt{bounds/} --- several shapefiles representing geographic areas
  of interest within the Lake Erie region
\end{itemize}

\textbf{Before writing any code, read the metadata file at}

\textbf{\texttt{../data/wildcard\_thursday/erie\_raster\_ex/erie\_README.txt}.}

This is not optional. Understanding your data source, its scaling, and
its data key is essential for correct analysis.

\hypertarget{introduction}{%
\subsection{Introduction}\label{introduction}}

Cyanobacteria (blue-green algae) blooms in Lake Erie have become a
significant environmental and public health concern. These blooms can
produce toxins that contaminate drinking water and harm aquatic
ecosystems. The data you will use were produced by NOAA using imagery
from the European Space Agency's Sentinel-3 satellite. Each raster
represents the maximum observed bloom intensity during a given week.

The product you will be working with is the \textbf{CIcyano}
(Cyanobacteria Index, cyano only) product. Values in these rasters are
stored as \textbf{digital numbers (DN)} and must be transformed to
meaningful physical units before interpretation. You saw this in class
--- recall the reverse-scaling formula from the README and from the Week
12 notebook.

A critical aspect of working with remote sensing data is understanding
\textbf{sentinel values}: special numeric codes that indicate conditions
like cloud cover, land pixels, or missing data. These are NOT real
observations and must be handled carefully before any analysis.

We begin by loading the relevant packages:

\begin{Shaded}
\begin{Highlighting}[]
\FunctionTok{library}\NormalTok{(terra)}
\end{Highlighting}
\end{Shaded}

\begin{verbatim}
## terra 1.8.60
\end{verbatim}

\begin{Shaded}
\begin{Highlighting}[]
\FunctionTok{library}\NormalTok{(sf)}
\end{Highlighting}
\end{Shaded}

\begin{verbatim}
## Linking to GEOS 3.13.0, GDAL 3.8.5, PROJ 9.5.1; sf_use_s2() is TRUE
\end{verbatim}

\begin{Shaded}
\begin{Highlighting}[]
\FunctionTok{library}\NormalTok{(tidyverse)}
\end{Highlighting}
\end{Shaded}

\begin{verbatim}
## -- Attaching core tidyverse packages ------------------------ tidyverse 2.0.0 --
## v dplyr     1.1.4     v readr     2.1.5
## v forcats   1.0.0     v stringr   1.5.1
## v ggplot2   3.5.2     v tibble    3.2.1
## v lubridate 1.9.3     v tidyr     1.3.1
## v purrr     1.0.2
\end{verbatim}

\begin{verbatim}
## -- Conflicts ------------------------------------------ tidyverse_conflicts() --
## x tidyr::extract() masks terra::extract()
## x dplyr::filter()  masks stats::filter()
## x dplyr::lag()     masks stats::lag()
## i Use the conflicted package (<http://conflicted.r-lib.org/>) to force all conflicts to become errors
\end{verbatim}

\begin{Shaded}
\begin{Highlighting}[]
\FunctionTok{library}\NormalTok{(tmap)}
\FunctionTok{library}\NormalTok{(fs)}
\end{Highlighting}
\end{Shaded}

\begin{center}\rule{0.5\linewidth}{0.5pt}\end{center}

\hypertarget{task-1-data-management-and-preprocessing}{%
\subsection{Task 1: Data management and
preprocessing}\label{task-1-data-management-and-preprocessing}}

\hypertarget{inventory-your-rasters}{%
\subsubsection{1.1 Inventory your
rasters}\label{inventory-your-rasters}}

Use \texttt{fs::dir\_ls} with the parameter
\texttt{regexp\ =\ "\textbackslash{}\textbackslash{}.tif\$"} to get the
file paths for all raster images in the \texttt{images/} directory.
Store the result in an object called \texttt{raster\_files}. Print the
result --- how many files are there?

\begin{Shaded}
\begin{Highlighting}[]
\CommentTok{\# your code here}

\NormalTok{raster\_files }\OtherTok{\textless{}{-}}\NormalTok{ fs}\SpecialCharTok{::}\FunctionTok{dir\_ls}\NormalTok{(}\StringTok{"erie\_raster\_ex/images/"}\NormalTok{, }\AttributeTok{regexp =} \StringTok{"}\SpecialCharTok{\textbackslash{}\textbackslash{}}\StringTok{.tif$"}\NormalTok{)}

\CommentTok{\# print file paths}
\NormalTok{raster\_files}
\end{Highlighting}
\end{Shaded}

\begin{verbatim}
## erie_raster_ex/images/ts_2017.0802_0808.L4.LE3.CIcyano.MAXIMUM_7day.tif
## erie_raster_ex/images/ts_2017.0816_0822.L4.LE3.CIcyano.MAXIMUM_7day.tif
## erie_raster_ex/images/ts_2017.0823_0829.L4.LE3.CIcyano.MAXIMUM_7day.tif
## erie_raster_ex/images/ts_2017.0920_0926.L4.LE3.CIcyano.MAXIMUM_7day.tif
## erie_raster_ex/images/ts_2017.0927_1003.L4.LE3.CIcyano.MAXIMUM_7day.tif
## erie_raster_ex/images/ts_2017.1018_1024.L4.LE3.CIcyano.MAXIMUM_7day.tif
## erie_raster_ex/images/ts_2017.1025_1031.L4.LE3.CIcyano.MAXIMUM_7day.tif
\end{verbatim}

\begin{Shaded}
\begin{Highlighting}[]
\CommentTok{\# number of files}
\FunctionTok{length}\NormalTok{(raster\_files)}
\end{Highlighting}
\end{Shaded}

\begin{verbatim}
## [1] 7
\end{verbatim}

\hypertarget{inspect-a-single-raster}{%
\subsubsection{1.2 Inspect a single
raster}\label{inspect-a-single-raster}}

Load the first raster in your list using \texttt{terra::rast()}. Examine
the following properties and record what you find:

\begin{itemize}
\tightlist
\item
  Coordinate reference system: \texttt{terra::crs()}
\item
  Spatial extent: \texttt{terra::ext()}
\item
  Spatial resolution: \texttt{terra::res()}
\item
  Number of layers: \texttt{terra::nlyr()}
\item
  Range of cell values (use \texttt{terra::global()} with \texttt{"min"}
  and \texttt{"max"})
\end{itemize}

\begin{Shaded}
\begin{Highlighting}[]
\CommentTok{\# your code here}
\CommentTok{\# load first raster}
\NormalTok{r1 }\OtherTok{\textless{}{-}}\NormalTok{ terra}\SpecialCharTok{::}\FunctionTok{rast}\NormalTok{(raster\_files[}\DecValTok{1}\NormalTok{])}

\CommentTok{\# Coords}
\NormalTok{terra}\SpecialCharTok{::}\FunctionTok{crs}\NormalTok{(r1)}
\end{Highlighting}
\end{Shaded}

\begin{verbatim}
## [1] "PROJCRS[\"WGS 84 / UTM zone 17N\",\n    BASEGEOGCRS[\"WGS 84\",\n        DATUM[\"World Geodetic System 1984\",\n            ELLIPSOID[\"WGS 84\",6378137,298.257223563,\n                LENGTHUNIT[\"metre\",1]]],\n        PRIMEM[\"Greenwich\",0,\n            ANGLEUNIT[\"degree\",0.0174532925199433]],\n        ID[\"EPSG\",4326]],\n    CONVERSION[\"Transverse Mercator\",\n        METHOD[\"Transverse Mercator\",\n            ID[\"EPSG\",9807]],\n        PARAMETER[\"Latitude of natural origin\",0,\n            ANGLEUNIT[\"degree\",0.0174532925199433],\n            ID[\"EPSG\",8801]],\n        PARAMETER[\"Longitude of natural origin\",-81,\n            ANGLEUNIT[\"degree\",0.0174532925199433],\n            ID[\"EPSG\",8802]],\n        PARAMETER[\"Scale factor at natural origin\",0.9996,\n            SCALEUNIT[\"unity\",1],\n            ID[\"EPSG\",8805]],\n        PARAMETER[\"False easting\",500000,\n            LENGTHUNIT[\"metre\",1],\n            ID[\"EPSG\",8806]],\n        PARAMETER[\"False northing\",0,\n            LENGTHUNIT[\"metre\",1],\n            ID[\"EPSG\",8807]]],\n    CS[Cartesian,2],\n        AXIS[\"easting\",east,\n            ORDER[1],\n            LENGTHUNIT[\"metre\",1]],\n        AXIS[\"northing\",north,\n            ORDER[2],\n            LENGTHUNIT[\"metre\",1]],\n    ID[\"EPSG\",32617]]"
\end{verbatim}

\begin{Shaded}
\begin{Highlighting}[]
\CommentTok{\# extent}
\NormalTok{terra}\SpecialCharTok{::}\FunctionTok{ext}\NormalTok{(r1)}
\end{Highlighting}
\end{Shaded}

\begin{verbatim}
## SpatExtent : 281282.934996, 694382.934996, 4572608.997997, 4764308.997997 (xmin, xmax, ymin, ymax)
\end{verbatim}

\begin{Shaded}
\begin{Highlighting}[]
\CommentTok{\# resolution}
\NormalTok{terra}\SpecialCharTok{::}\FunctionTok{res}\NormalTok{(r1)}
\end{Highlighting}
\end{Shaded}

\begin{verbatim}
## [1] 300 300
\end{verbatim}

\begin{Shaded}
\begin{Highlighting}[]
\CommentTok{\# layers}
\NormalTok{terra}\SpecialCharTok{::}\FunctionTok{nlyr}\NormalTok{(r1)}
\end{Highlighting}
\end{Shaded}

\begin{verbatim}
## [1] 1
\end{verbatim}

\begin{Shaded}
\begin{Highlighting}[]
\CommentTok{\# min and max}
\NormalTok{terra}\SpecialCharTok{::}\FunctionTok{global}\NormalTok{(r1, }\FunctionTok{c}\NormalTok{(}\StringTok{"min"}\NormalTok{, }\StringTok{"max"}\NormalTok{), }\AttributeTok{na.rm =} \ConstantTok{TRUE}\NormalTok{)}
\end{Highlighting}
\end{Shaded}

\begin{verbatim}
##                                               min max
## ts_2017.0802_0808.L4.LE3.CIcyano.MAXIMUM_7day   1 252
\end{verbatim}

\hypertarget{write-a-preprocessing-function-clean_cicyano}{%
\subsubsection{\texorpdfstring{1.3 Write a preprocessing function:
\texttt{clean\_cicyano()}}{1.3 Write a preprocessing function: clean\_cicyano()}}\label{write-a-preprocessing-function-clean_cicyano}}

Looking at the data key in the README, you will notice that values of 0
and 250 and above have special meanings --- they are \textbf{not} valid
bloom observations. Before doing any analysis, these sentinel values
must be replaced with \texttt{NA}.

Write a function called \texttt{clean\_cicyano(r)} that:

\begin{enumerate}
\def\labelenumi{\arabic{enumi}.}
\tightlist
\item
  Takes a SpatRaster \texttt{r} as input
\item
  Builds a reclassification matrix that maps \textbf{0} and \textbf{all
  values \textgreater= 250} to \texttt{NA}, and leaves values 1--249
  unchanged
\item
  Applies the reclassification using \texttt{terra::classify()} with the
  argument \texttt{others\ =\ NA} (this is a useful argument --- look it
  up)
\item
  Returns the cleaned raster
\end{enumerate}

\emph{Hint: Your reclassification matrix should handle the no-detection
case (0) and all sentinel codes (250 through 255). Review how
\texttt{terra::classify()} uses a three-column matrix to specify
\texttt{{[}from,\ to,\ becomes{]}} intervals.}

\begin{Shaded}
\begin{Highlighting}[]
\CommentTok{\# your code here}
\NormalTok{clean\_cicyano }\OtherTok{\textless{}{-}} \ControlFlowTok{function}\NormalTok{(r) \{}
\NormalTok{  rcl }\OtherTok{\textless{}{-}} \FunctionTok{matrix}\NormalTok{(}\FunctionTok{c}\NormalTok{(}
    \DecValTok{0}\NormalTok{, }\DecValTok{0}\NormalTok{, }\ConstantTok{NA}\NormalTok{,}
    \DecValTok{249}\NormalTok{, }\DecValTok{255}\NormalTok{, }\ConstantTok{NA}
\NormalTok{  ), }\AttributeTok{ncol =} \DecValTok{3}\NormalTok{, }\AttributeTok{byrow =} \ConstantTok{TRUE}\NormalTok{)}
\NormalTok{  r\_clean }\OtherTok{\textless{}{-}}\NormalTok{ terra}\SpecialCharTok{::}\FunctionTok{classify}\NormalTok{(r, rcl)}
  \FunctionTok{return}\NormalTok{(r\_clean)}
\NormalTok{\}}


\CommentTok{\# test}
\NormalTok{r1\_clean }\OtherTok{\textless{}{-}} \FunctionTok{clean\_cicyano}\NormalTok{(r1)}

\NormalTok{terra}\SpecialCharTok{::}\FunctionTok{global}\NormalTok{(r1, }\FunctionTok{c}\NormalTok{(}\StringTok{"min"}\NormalTok{, }\StringTok{"max"}\NormalTok{), }\AttributeTok{na.rm =} \ConstantTok{TRUE}\NormalTok{)}
\end{Highlighting}
\end{Shaded}

\begin{verbatim}
##                                               min max
## ts_2017.0802_0808.L4.LE3.CIcyano.MAXIMUM_7day   1 252
\end{verbatim}

\begin{Shaded}
\begin{Highlighting}[]
\NormalTok{terra}\SpecialCharTok{::}\FunctionTok{global}\NormalTok{(r1\_clean, }\FunctionTok{c}\NormalTok{(}\StringTok{"min"}\NormalTok{, }\StringTok{"max"}\NormalTok{), }\AttributeTok{na.rm =} \ConstantTok{TRUE}\NormalTok{)}
\end{Highlighting}
\end{Shaded}

\begin{verbatim}
##                                               min max
## ts_2017.0802_0808.L4.LE3.CIcyano.MAXIMUM_7day   1 249
\end{verbatim}

Test your function on the first raster. Verify it worked by comparing
the global min/max of the raw raster vs.~the cleaned raster.

\hypertarget{write-a-transform-function-dn_to_ci}{%
\subsubsection{\texorpdfstring{1.4 Write a transform function:
\texttt{dn\_to\_ci()}}{1.4 Write a transform function: dn\_to\_ci()}}\label{write-a-transform-function-dn_to_ci}}

The DN values in these rasters are scaled according to the formula in
the README. To convert back to physical CI units, use the
reverse-scaling formula:

\[CI = 10^{\left(\frac{3.0}{250.0} \times DN - 4.2\right)}\]

Write a function called \texttt{dn\_to\_ci(r)} that takes a cleaned
SpatRaster and returns a new raster with values transformed to CI units.
You can apply a function to a raster using standard arithmetic
operations or by passing a function to the raster.

\begin{Shaded}
\begin{Highlighting}[]
\CommentTok{\# your code here}

\NormalTok{dn\_to\_ci }\OtherTok{\textless{}{-}} \ControlFlowTok{function}\NormalTok{(r) \{}
\NormalTok{  ci }\OtherTok{\textless{}{-}} \DecValTok{10}\SpecialCharTok{\^{}}\NormalTok{((}\FloatTok{3.0} \SpecialCharTok{/} \FloatTok{250.0}\NormalTok{) }\SpecialCharTok{*}\NormalTok{ r }\SpecialCharTok{{-}} \FloatTok{4.2}\NormalTok{)}
  \FunctionTok{return}\NormalTok{(ci)}
\NormalTok{\}}

\CommentTok{\# test}
\NormalTok{r1\_ci }\OtherTok{\textless{}{-}} \FunctionTok{dn\_to\_ci}\NormalTok{(r1\_clean)}

\NormalTok{terra}\SpecialCharTok{::}\FunctionTok{global}\NormalTok{(r1\_ci, }\FunctionTok{c}\NormalTok{(}\StringTok{"min"}\NormalTok{, }\StringTok{"max"}\NormalTok{), }\AttributeTok{na.rm =} \ConstantTok{TRUE}\NormalTok{)}
\end{Highlighting}
\end{Shaded}

\begin{verbatim}
##                                                        min       max
## ts_2017.0802_0808.L4.LE3.CIcyano.MAXIMUM_7day 6.486344e-05 0.0613762
\end{verbatim}

Test your function. What is the range of CI values in the transformed
raster? Do they make physical sense given the data key?

\hypertarget{questions-for-task-1}{%
\subsubsection{Questions for Task 1:}\label{questions-for-task-1}}

\textbf{Q1.1} What does the CRS of these rasters tell you about how
distances and areas are measured? Would you need to reproject before
calculating bloom area in square kilometers?

\textbf{Q1.2} Why is it important to remove sentinel values
\emph{before} applying the DN-to-CI transform --- rather than after?
What would happen to a land pixel (value = 252) if you transformed it
without cleaning first?

\textbf{Q1.3} The \texttt{others\ =\ NA} argument in
\texttt{terra::classify()} is useful here. In your own words, what does
it do?

\begin{center}\rule{0.5\linewidth}{0.5pt}\end{center}

\hypertarget{task-2-spatial-subsetting-crop-and-mask}{%
\subsection{Task 2: Spatial subsetting --- crop and
mask}\label{task-2-spatial-subsetting-crop-and-mask}}

\hypertarget{load-and-align-an-area-of-interest}{%
\subsubsection{2.1 Load and align an area of
interest}\label{load-and-align-an-area-of-interest}}

Load the \texttt{bb\_erie\_sandusky} shapefile from the \texttt{bounds/}
directory using \texttt{sf::read\_sf()}. Check its CRS. If it does not
match the raster CRS, reproject it using \texttt{sf::st\_transform()}.

\begin{Shaded}
\begin{Highlighting}[]
\CommentTok{\# your code here}
\end{Highlighting}
\end{Shaded}

\hypertarget{crop-and-mask}{%
\subsubsection{2.2 Crop and mask}\label{crop-and-mask}}

Using the \textbf{August 16--22, 2017} raster
(\texttt{ts\_2017.0816\_0822}):

\begin{enumerate}
\def\labelenumi{\arabic{enumi}.}
\tightlist
\item
  Apply \texttt{clean\_cicyano()} and \texttt{dn\_to\_ci()} to get a
  clean, transformed raster
\item
  Use \texttt{terra::crop()} to clip the raster extent to the Sandusky
  Bay boundary
\item
  Use \texttt{terra::mask()} to set pixels \emph{outside} the boundary
  to \texttt{NA}
\end{enumerate}

Store your final result as \texttt{sandusky\_0816}.

\begin{Shaded}
\begin{Highlighting}[]
\CommentTok{\# your code here}
\end{Highlighting}
\end{Shaded}

\hypertarget{compare-full-vs.-masked}{%
\subsubsection{2.3 Compare full
vs.~masked}\label{compare-full-vs.-masked}}

Using \texttt{terra::global()}, calculate and compare the \textbf{mean}
CI value for:

\begin{itemize}
\tightlist
\item
  The full Lake Erie cleaned/transformed raster
\item
  The Sandusky Bay-masked raster (\texttt{sandusky\_0816})
\end{itemize}

\begin{Shaded}
\begin{Highlighting}[]
\CommentTok{\# your code here}
\end{Highlighting}
\end{Shaded}

\hypertarget{map-it}{%
\subsubsection{2.4 Map it}\label{map-it}}

Create a \texttt{tmap} map of \texttt{sandusky\_0816}. Your map should
include:

\begin{itemize}
\tightlist
\item
  A meaningful title
\item
  A legend with appropriate labels
\item
  A scale bar
\end{itemize}

\begin{Shaded}
\begin{Highlighting}[]
\CommentTok{\# your code here}
\end{Highlighting}
\end{Shaded}

\hypertarget{questions-for-task-2}{%
\subsubsection{Questions for Task 2:}\label{questions-for-task-2}}

\textbf{Q2.1} In your own words, what is the difference between
\texttt{terra::crop()} and \texttt{terra::mask()}? Why did we use both
here instead of just one?

\textbf{Q2.2} The mean CI value for Sandusky Bay is likely different
from the full lake mean. What could explain this? Does this necessarily
mean Sandusky Bay has more or less cyanobacteria than average?

\begin{center}\rule{0.5\linewidth}{0.5pt}\end{center}

\hypertarget{task-3-multi-temporal-comparison}{%
\subsection{Task 3: Multi-temporal
comparison}\label{task-3-multi-temporal-comparison}}

\hypertarget{load-and-stack-multiple-rasters}{%
\subsubsection{3.1 Load and stack multiple
rasters}\label{load-and-stack-multiple-rasters}}

Load and clean \textbf{at least four} of the seven rasters (choose dates
that span a meaningful time range). Apply \texttt{clean\_cicyano()} and
\texttt{dn\_to\_ci()} to each. Then combine them into a single
multi-layer SpatRaster using \texttt{c()}.

\emph{Tip: using \texttt{purrr::map()} over your \texttt{raster\_files}
vector, combined with your preprocessing functions, can make this more
efficient.}

\begin{Shaded}
\begin{Highlighting}[]
\CommentTok{\# your code here}
\end{Highlighting}
\end{Shaded}

\hypertarget{temporal-summary}{%
\subsubsection{3.2 Temporal summary}\label{temporal-summary}}

Use \texttt{terra::global()} on your multi-layer raster to calculate the
\textbf{mean} CI value per layer (use \texttt{na.rm\ =\ TRUE}). This
will give you one mean value per date.

Organize the results into a data frame with at minimum two columns:
\texttt{date} and \texttt{mean\_ci}. You will need to extract the date
from each filename. Use \texttt{basename()} and/or
\texttt{stringr::str\_extract()} to parse the date string from the file
path --- the date range is the part of the filename that looks like
\texttt{0802\_0808}.

\begin{Shaded}
\begin{Highlighting}[]
\CommentTok{\# your code here}
\end{Highlighting}
\end{Shaded}

\hypertarget{plot-the-time-series}{%
\subsubsection{3.3 Plot the time series}\label{plot-the-time-series}}

Using \texttt{ggplot2}, create a line or point plot showing mean CI
values over time. Label your axes clearly. Which date in your dataset
shows the highest mean bloom intensity?

\begin{Shaded}
\begin{Highlighting}[]
\CommentTok{\# your code here}
\end{Highlighting}
\end{Shaded}

\hypertarget{binary-bloom-comparison}{%
\subsubsection{3.4 Binary bloom
comparison}\label{binary-bloom-comparison}}

Choose any \textbf{two dates} from your dataset. For each, create a
binary raster where cells with valid CI values are coded as \texttt{1}
and all other cells (NA) as \texttt{0}. You will need to think about how
to do this using \texttt{terra::classify()} or logical raster
operations.

Use these binary rasters to create two new rasters:

\begin{itemize}
\tightlist
\item
  \textbf{Persistent bloom:} cells where blooms were present on
  \emph{both} dates
\item
  \textbf{New bloom:} cells where blooms were present on the
  \emph{later} date but NOT the earlier date
\end{itemize}

Map both results.

\emph{This requires careful logical thinking before writing code. Plan
your approach in pseudocode first.}

\begin{Shaded}
\begin{Highlighting}[]
\CommentTok{\# your code here}
\end{Highlighting}
\end{Shaded}

\hypertarget{questions-for-task-3}{%
\subsubsection{Questions for Task 3:}\label{questions-for-task-3}}

\textbf{Q3.1} What are the limitations of using \textbf{mean CI value}
as a measure of bloom severity across the whole lake? What would be a
more robust alternative metric?

\textbf{Q3.2} Stacking rasters with \texttt{c()} in terra creates a
multi-layer SpatRaster. What requirement must all rasters meet in order
to be stacked this way? (Think: extent, resolution, CRS.)

\begin{center}\rule{0.5\linewidth}{0.5pt}\end{center}

\hypertarget{task-4-zonal-analysis-with-areas-of-interest}{%
\subsection{Task 4: Zonal analysis with areas of
interest}\label{task-4-zonal-analysis-with-areas-of-interest}}

\hypertarget{load-all-boundary-files}{%
\subsubsection{4.1 --- Load all boundary
files}\label{load-all-boundary-files}}

Use \texttt{list.files()} to get all \texttt{.shp} files from the
\texttt{bounds/} directory. Load each with \texttt{sf::read\_sf()}, then
combine them into a single sf object using \texttt{dplyr::bind\_rows()}.
Make sure each polygon has an identifier --- if not, add a column with
the filename or AOI name.

\begin{Shaded}
\begin{Highlighting}[]
\CommentTok{\# your code here}
\end{Highlighting}
\end{Shaded}

\hypertarget{extract-by-area}{%
\subsubsection{4.2 Extract by area}\label{extract-by-area}}

Using the date you identified in Task 3.3 as having the highest mean
bloom intensity, extract CI values for each area of interest using
\texttt{terra::extract()}. Calculate the \textbf{mean} and
\textbf{maximum} CI per AOI.

\emph{Tip: \texttt{terra::extract()} returns a data frame. The
\texttt{ID} column links rows to the AOI that was used for extraction.
You can use \texttt{group\_by()} and \texttt{summarize()} to aggregate.}

\begin{Shaded}
\begin{Highlighting}[]
\CommentTok{\# your code here}
\end{Highlighting}
\end{Shaded}

\hypertarget{summary-table}{%
\subsubsection{4.3 Summary table}\label{summary-table}}

Present your results as a clean summary table. Which AOI had the highest
mean CI value on that date? Which had the highest maximum?

\begin{Shaded}
\begin{Highlighting}[]
\CommentTok{\# your code here}
\end{Highlighting}
\end{Shaded}

\hypertarget{map-with-aoi-overlays}{%
\subsubsection{4.4 Map with AOI overlays}\label{map-with-aoi-overlays}}

Create a \texttt{tmap} map showing the CI raster for your chosen date
with the AOI boundaries overlaid. Use \texttt{tm\_borders()} for the
boundaries and label them if possible.

\begin{Shaded}
\begin{Highlighting}[]
\CommentTok{\# your code here}
\end{Highlighting}
\end{Shaded}

\hypertarget{questions-for-task-4}{%
\subsubsection{Questions for Task 4:}\label{questions-for-task-4}}

\textbf{Q4.1} You may have gotten \texttt{NA} values in your
\texttt{terra::extract()} results for some AOIs. What could cause this?
Name at least two reasons.

\textbf{Q4.2} Looking at the spatial distribution of your AOIs on the
map, which areas do you expect to be most prone to severe blooms, and
why? (Consider geography, hydrology, land use in the Lake Erie
watershed.)

\begin{center}\rule{0.5\linewidth}{0.5pt}\end{center}

\hypertarget{questions}{%
\subsection{Questions}\label{questions}}

\begin{enumerate}
\def\labelenumi{\arabic{enumi}.}
\item
  Throughout this lab you have been working with \textbf{raster} data.
  In what ways was the analysis workflow different from working with
  vector data (like you did in Labs 1--3)? What concepts translated
  directly, and what was fundamentally different?
\item
  Cyanobacteria blooms are affected by temperature, nutrient runoff, and
  lake circulation. Based on the temporal patterns you observed in Task
  3, does the timing of peak bloom match what you would expect
  ecologically? Explain your reasoning (I know you're not an expert on
  blooms - just demonstrate thinking).
\item
  We used \texttt{terra::extract()} in Task 4 to perform a zonal
  summary. How does this relate conceptually to the \textbf{zonal
  operations} category of map algebra discussed in class?
\end{enumerate}

\end{document}
